\documentclass[10pt,a4paper]{amsart}
\usepackage[margin=19mm]{geometry}
\usepackage{amsmath,amssymb,mathtools}
\usepackage[scr=rsfs,cal=euler]{mathalfa}
\usepackage{xfrac}
\usepackage[unicode,hidelinks,bookmarks=false]{hyperref}
\newcommand{\ie}{i.e.\ }
\newcommand{\cf}{cf.\ }
\newcommand{\resp}{respectively\ }
\newcommand{\Subsection}{\subsection}
\providecommand{\nobreakdash}{\nobreak\hbox{-}}
\setlength{\emergencystretch}{2em}
\multlinegap0pt
\usepackage{tikz}
\usetikzlibrary{cd}
\usepackage{mathtools}
\usepackage{amssymb}
\usepackage{xcolor}
\newtheorem{definition}{Definition}
\newtheorem{proposition}{Proposition}
\DeclareMathOperator{\Bdrp}{B^+_{dR}}
\DeclareMathOperator{\C}{\textbf{C}}
\newcommand{\mBdrp}{B^+_{\mathrm{dR},\Delta}}
\newcommand{\pmBdrp}{B^{+,(p)}_{\mathrm{dR},\Delta}}
\title{Multivariate Period Rings}
\author{Rohit Pokhrel}
\date{Source: arXiv:2606.11706v1 (CC BY 4.0)}
\begin{document}
\maketitle

\noindent\emph{Source and licence.} Multivariate Period Rings. Version arXiv:2606.11706v1 (CC BY 4.0), distributed by arXiv under the Creative Commons Attribution 4.0 International licence, http://creativecommons.org/licenses/by/4.0/, as recorded in the arXiv licence metadata for this exact version and shown on the abstract page https://arxiv.org/abs/2606.11706v1. Full licence notice: \texttt{licenses/arxiv-2606.11706-CC-BY-4.0.txt}.\par\smallskip

\noindent\textit{Editorial scope.} Counted unit: the article's proposition proposition (source0.tex lines 520--522). The article's complete original proof of it is delivered with it (source0.tex lines 524--548, one \texttt{proof} environment of 25 lines). No mathematics is written, repaired or re-derived: the statement and the proof are the article's own lines, in the article's own notation, and no retained line is substituted or re-flowed.  Retained uncounted as the source-local context the proof needs: lines 481--486.  Nothing was omitted from any retained span, so the package carries no omission record.  No retained line contains a citation, so the wrapper carries no bibliography and no bibliography entry.  The retained lines contain the article's own diagram code, which the wrapper typesets with the standard tikz packages; no image file is delivered and the package declares no asset dependency.  Editorial formatting only, in addition: an AMS \texttt{amsart} preamble that restates the article's own declarations and macros used by the retained lines, namely the environments \texttt{definition}, \texttt{proposition} on separate counters, together with the standard packages those lines need.  The retained environments print in this document as Definition 1, Proposition 1 in order of appearance, whereas the article prints the same items with the numbering of its own full document.  The complete native source of the article as distributed in the arXiv e-print for this exact version is bundled unchanged as \texttt{sources/202569/source0.tex}.
\begin{definition}\label{Profinite completion}
    Let $R$ be a ring and let $\{\mathfrak{m}_i\}_{i\in\Lambda}$ be a family of maximal ideals of $R$ such that $\bigcap\limits_{i\in\Lambda}\mathfrak{m}_\lambda=\{0\}$. Let $\Lambda'\subset\Lambda$ be a finite subset, and define $\mathfrak{a}_{\Lambda'}=\bigcap\limits_{i\in\Lambda'}\mathfrak{m}_i$, then the \emph{profinite completion} of $R$ with respect to the family $\{\mathfrak{m}_i\}_{i\in\Lambda}$ is defined to be
    \begin{align*}
        \hat{R}:=\varprojlim\limits_{\Lambda'\subset\Lambda\text{ finite }}R/\mathfrak{a}_{\Lambda'}.
    \end{align*}
\end{definition}

\begin{proposition}
    The map $\theta_\Delta:\pmBdrp\to\C_\Delta$ has an extension to a surjective ring homomorphism $\theta_\Delta:\mBdrp\to\widehat{\C}_\Delta$, where $\widehat{\C}_\Delta$ is the profinite completion of $\C_\Delta$ with respect to the family $\{\mathfrak{P}_s\}_{s\in\mathfrak{S}}$ in sense of Definition \ref{Profinite completion}.
\end{proposition}

\begin{proof}
    We have the commutative diagram
    \[
    \begin{tikzcd}
        \pmBdrp\arrow[r,"\theta_\Delta"]\arrow[d,"\Phi_s"]& \C_\Delta\arrow[d,"\beta_s"]\\
        \Bdrp\arrow[r,"\theta"]  &\C
    \end{tikzcd}
    \]
    i.e., $\beta_s\circ\theta_\Delta=\theta\circ\Phi_s$. Then we have
    \begin{align*}
        \theta^{-1}_\Delta(\mathfrak{P}_s)=\ker(\beta_s\circ\theta_\Delta)=\ker(\theta\circ\Phi_s)=\Phi^{-1}_s(\ker\theta)\supset \ker\Phi_s.
    \end{align*}
    Then for any finite subset $\mathfrak{S}'$ of $\mathfrak{S}$, define a ring homomorphism
    \begin{align*}
        \pmBdrp/\bigcap\limits_{s\in\mathfrak{S}'}\ker\Phi_s&\to \C_\Delta/\bigcap\limits_{s\in\mathfrak{S}'}\mathfrak{P}_s,\\
        x+\bigcap\limits_{s\in\mathfrak{S}'}\ker\Phi_s&\mapsto \theta_\Delta(x)+\bigcap\limits_{s\in\mathfrak{S}'}\mathfrak{P}_s.
    \end{align*}
    We want to show that this map is well-defined. For this, suppose that $x+\bigcap\limits_{s\in\mathfrak{S}'}\ker\Phi_s=y+\bigcap\limits_{s\in\mathfrak{S}'}\ker\Phi_s$. Then we have
    \begin{align*}
        x-y\in\ker\Phi_s\text{ for all }s\in\mathfrak{S}'&\implies x-y\in\ker\Phi_s\text{ for all }s\in\mathfrak{S}'\\
        &\implies x-y\in\theta^{-1}_\Delta(\mathfrak{P}_s)\text{ for all }s\in\mathfrak{S}'\\
        &\implies \theta_\Delta(x)-\theta_\Delta(y)\in\mathfrak{P}_s\text{ for all }s\in\mathfrak{S'}.
    \end{align*}
    This shows that $\theta_\Delta(x)+\bigcap\limits_{s\in\mathfrak{S}'}\mathfrak{P}_s=\theta_\Delta(y)+\bigcap\limits_{s\in\mathfrak{S}'}\mathfrak{P}_s$. This shows that the above map is well-defined, and it is clear that this extends the ring map $\theta_\Delta:\pmBdrp\to\C_\Delta$, and taking the inverse limit we get the required ring homomorphism. Also note that $\ker\theta_\Delta$ is principal, which is generated by the system of images of $\xi$ onto each quotient, which we also denote by $\xi$. We denote this map also by $\theta_\Delta$.
\end{proof}

\end{document}
